\documentclass[10pt,a4paper]{article}
\usepackage[utf8]{inputenc}
\usepackage{geometry}
\usepackage{times} % Using Times Roman font
\usepackage{titling}

% Set the page size and margins
\geometry{
  a4paper,
  top=30mm,
  bottom=30mm,
  left=30mm,
  right=30mm
}

% Set title formatting
\pretitle{\begin{center}\large\bfseries}
\posttitle{\par\end{center}\vskip 0.5em}
\preauthor{\begin{center}\small}
\postauthor{\par\end{center}}
\predate{\par\begin{center}\small}
\postdate{\par\end{center}}

\title{Title in Times Roman 14 point – Upper and Lower Case}
\author{Presenting Author, Co-Authors\\
\textit{Affiliation}\\
\textit{Address}\\
\textit{e-mail address}}
\date{} % No date

\begin{document}

\maketitle

\begin{abstract}

oindent These instructions are an example of what a properly prepared meeting abstract should look like. Proper column and margin measurements are indicated.

The abstract should not exceed ONE PAGE of text, references, tables, and figures. Abstracts exceeding this limit may be cut without consideration of content after the first page.

\textbf{Paper Size:} A4 (21.0 x 29.7 cm)

\textbf{Margins} \\
Top: 30.0 mm \\
Bottom: 30.0 mm \\
Sides: 30.0 mm
\end{abstract}


oindent Type the body of the abstract text (including references and tables) single-spaced in 10-point Times Roman regular. The abstract should be concise, clear, and based on the most recent research related to the topic. It should contain enough information for a reader to understand the context of the research, the methods used, the main findings, and the significance of the conclusions drawn from the study.

\end{document}
